\section{\\ Исследование архитектуры пикосети Bluetooth и организация обмена \\данными через RFCOMM}

\subsection{Цель работы}
Изучить принципы построения Bluetooth-пикосетей, технологию скачкообразной перестройки рабочей частоты (FHSS) и механизмы установления логических соединений. Получить практические навыки программирования сетевого взаимодействия в архитектуре «мастер–ведомые» с использованием Python и стека BlueZ.

\subsection{Задачи}
\begin{itemize}[nosep]
    \item Изучить топологию пикосети, роли устройств (мастер/слейв) и ограничения архитектуры.
    \item Рассмотреть принцип работы FHSS, структуру частотной сетки и синхронизацию каналов.
    \item Ознакомиться с протоколом RFCOMM, его назначением и особенностями эмуляции последовательного порта.
    \item Настроить два модуля ESP32/ESP8266 в режиме Bluetooth Classic SPP-серверов.
    \item Разработать и запустить Python-скрипт на ПК, устанавливающий параллельные RFCOMM-соединения с обоими устройствами.
    \item Организовать двунаправленный обмен данными, протоколировать процесс и проанализировать надёжность канала.
\end{itemize}

\subsection{Теоретический материал}

\subsubsection{Архитектура пикосети и роли устройств}
Базовой единицей любой сети Bluetooth является \textbf{пикосеть} (piconet). Она представляет собой логическую сеть, образованную одним ведущим устройством (\textbf{мастером}) и до семи активных ведомых устройств (\textbf{слейвов}). Дополнительно в пикосети могут находиться до 255 устройств в режиме ожидания (паркованные устройства, parked), не участвующие в активном обмене данными, но готовых к активации.

Мастер выполняет функции синхронизации: его внутренние часы и 48-битный Bluetooth-адрес (BD\_ADDR) используются для генерации псевдослучайной частотной последовательности и определения временных слотов. Все слейвы синхронизируются по часам мастера и следуют единому расписанию скачков частоты. В рамках одной пикосети обмен данными возможен только между мастером и слейвами; прямое взаимодействие слейв–слейв не поддерживается на физическом уровне.

\subsubsection{Технология FHSS и организация физического канала}
Для обеспечения помехоустойчивости и безопасности в сетях Bluetooth BR/EDR применяется технология \textbf{FHSS (Frequency Hopping Spread Spectrum)}. Рабочий диапазон 2,4–2,4835 ГГц делится на 79 каналов с шагом 1 МГц (в некоторых регионах — 23 или 25 каналов). Скорость скачков составляет 1600 переходов в секунду (в базовом режиме BR).

Последовательность каналов вычисляется по алгоритму, использующему 28-битное значение адреса мастера, 28-битные внутренние часы и 1-битный идентификатор режима. Таким образом, каждое устройство в пикосети «знает», на какой частоте будет находиться мастер в следующий временной слот. Параметры физического канала представлены в табл. 2.1.

\begin{table}[htbp]
\centering
\caption{Параметры физического канала и технологии FHSS}
\label{tab:fhss_params}
\begin{tabularx}{\textwidth}{lX}
\toprule
\textbf{Параметр} & \textbf{Значение / Описание} \\
\midrule
Частотный диапазон & 2400–2483,5 МГц (ISM-диапазон) \\
Количество каналов & 79 (шаг 1 МГц) \\
Скорость скачков & 1600 гоп/с (1,6 кГц) \\
Время удержания частоты & 625 мкс (1 слот) \\
Синхронизация & По часам мастера и его BD\_ADDR \\
Модуляция & GFSK ( Gaussian Frequency-Shift Keying ) \\
\bottomrule
\end{tabularx}
\end{table}

Организация канала на уровне контроллера включает: \textit{физический канал} (определяется частотной сеткой), \textit{логический транспортный канал} (асинхронный ACL или синхронный SCO), \textit{логический канал} (для конкретного соединения) и \textit{логический транспортный канал} (каналы данных и сигнализации). Мастер управляет распределением временных слотов между слейвами, используя пакетную передачу данных.

\subsubsection{Протокол RFCOMM и программная реализация}
Для передачи потоковых данных в архитектуре Bluetooth используется подуровень \textbf{RFCOMM}, эмулирующий последовательный порт (эмуляция RS-232). RFCOMM работает поверх L2CAP и использует концепцию \textit{Data Link Connection Identifiers} (DLCI) для мультиплексирования до 60 виртуальных последовательных соединений в рамках одного физического канала.

В ОС семейства Linux реализация RFCOMM доступна через системные сокеты. В языке Python для работы с Bluetooth Classic и установления RFCOMM-соединений используется стандартная библиотека \texttt{socket} с доменом \texttt{AF\_BLUETOOTH} и протоколом \texttt{BTPROTO\_RFCOMM}. После установления соединения программный интерфейс полностью идентичен работе с обычными TCP-сокетами, что упрощает разработку клиент-серверных приложений для обмена текстовыми и бинарными данными.

\subsection{Порядок выполнения лабораторной работы}

\textbf{Примечание.} Для выполнения работы требуется ПК с установленным Linux и стеком BlueZ, а также два микроконтроллера ESP32 с поддержкой Bluetooth Classic. Перед началом работы устройства должны быть сопряжены с ПК (см. Лабораторную работу 1).

\begin{enumerate}[nosep, leftmargin=*]
    \item Подготовьте микроконтроллеры ESP32. Прошейте их микропрограммой, реализующей профиль SPP (Serial Port Profile) в режиме сервера. Убедитесь, что модули находятся в режиме \texttt{discoverable} и \texttt{pairable}. Запишите 48-битные MAC-адреса обоих устройств (формат \texttt{XX:XX:XX:XX:XX:XX}).

    \item Установите необходимые системные пакеты на лабораторном ПК:
    \begin{lstlisting}
user@host:~$ sudo apt update
user@host:~$ sudo apt install python3 python3-pip libbluetooth-dev
user@host:~$ pip3 install --user pybluez
    \end{lstlisting}

    \item Создайте файл \texttt{bt\_master.py} и реализуйте скрипт, устанавливающий два параллельных RFCOMM-соединения. Базовая структура приведена в листинге 2.1. Замените \texttt{MAC\_ESP1} и \texttt{MAC\_ESP2} на реальные адреса ваших модулей.

    \item Запустите скрипт с правами суперпользователя (для доступа к HCI-интерфейсу):
    \begin{lstlisting}
user@host:~$ sudo python3 bt_master.py
    \end{lstlisting}

    \item Наблюдайте за выводом терминала. Скрипт должен последовательно подключиться к обоим устройствам, отправить тестовые пакеты и получить ответы. В случае разрыва соединения скрипт должен автоматически выполнить повторное подключение.

    \item В параллельном окне терминала запустите мониторинг трафика:
    \begin{lstlisting}
user@host:~$ sudo btmon -w lab02_trace.log
    \end{lstlisting}
    После завершения обмена остановите \texttt{btmon} и просмотрите лог, обращая внимание на этапы \texttt{L2CAP Connect Request}, \texttt{RFCOMM SABM} и \texttt{Encryption Change}.

    \item Зафиксируйте в отчёте успешные подключения, переданные/принятые байты и время отклика от каждого ведомого устройства.
\end{enumerate}

\begin{lstlisting}[caption={Базовая структура Python-скрипта для управления пикосетью}]
import socket
import time
import bluetooth

# MAC addresses of slave devices (replace with real ones)
ESP1_MAC = "AA:BB:CC:DD:EE:01"
ESP2_MAC = "AA:BB:CC:DD:EE:02"
RFCOMM_PORT = 1  # Standard port for SPP

def connect_rfcomm(mac_addr, port=RFCOMM_PORT):
    sock = socket.socket(socket.AF_BLUETOOTH, 
                         socket.SOCK_STREAM, 
                         socket.BTPROTO_RFCOMM)
    try:
        sock.connect((mac_addr, port))
        print(f"[+] Connected to {mac_addr}")
        return sock
    except Exception as e:
        print(f"[-] Connection error to {mac_addr}: {e}")
        return None

def exchange_data(sock, device_id):
    if not sock: return
    msg = f"Data from Master -> {device_id}\n"
    sock.send(msg.encode('utf-8'))
    response = sock.recv(1024).decode('utf-8').strip()
    print(f"<- Response from {device_id}: {response}")

def main():
    print("=== Piconet initialization ===")
    s1 = connect_rfcomm(ESP1_MAC)
    s2 = connect_rfcomm(ESP2_MAC)
    
    if s1 and s2:
        print("=== Both channels active. Starting exchange ===")
        for i in range(5):
            exchange_data(s1, "ESP32_1")
            exchange_data(s2, "ESP32_2")
            time.sleep(1)
        
        s1.close()
        s2.close()
        print("=== Connections closed ===")
    else:
        print("[-] Failed to establish all connections")

if __name__ == "__main__":
    main()
\end{lstlisting}

\subsection{Содержание отчета}
\begin{enumerate}[nosep, leftmargin=*]
    \item Заголовок согласно приложению.
    \item Цель работы.
    \item Задание согласно варианту.
    \item Листинги результатов работы:
    \begin{enumerate}[label=\alph*), nosep]
        \item схема топологии созданной пикосети с указанием ролей и MAC-адресов;
        \item листинг Python-скрипта с комментариями;
        \item вывод терминала с результатами подключения и обмена данными;
        \item фрагменты лога \texttt{btmon}, подтверждающие установление L2CAP- и RFCOMM-каналов, а также активацию шифрования;
        \item таблица времени отклика и потерь пакетов для каждого ведомого устройства.
    \end{enumerate}
    \item Ответы на контрольные вопросы.
    \item Краткие выводы по работе.
\end{enumerate}

\textbf{Примечание.} При оформлении отчёта допускается сокращать листинги до ключевых фрагментов. Вывод \texttt{btmon} рекомендуется прилагать в виде вырезок, иллюстрирующих этапы установки соединения и передачи данных.

\section*{Контрольные вопросы}
\begin{enumerate}[nosep, leftmargin=*]
    \item Что такое пикосеть и каковы её основные ограничения по количеству активных устройств?
    \item Какую роль выполняет мастер в организации пикосети? Чем его функции отличаются от функций слейва?
    \item Объясните принцип работы технологии FHSS. Почему скорость скачков в Bluetooth составляет именно 1600 гоп/с?
    \item Как адрес мастера и его внутренние часы влияют на генерацию частотной последовательности?
    \item Какую задачу решает протокол RFCOMM в стеке Bluetooth и как он взаимодействует с L2CAP?
    \item Почему при программировании сокетов в Python используется константа \texttt{BTPROTO\_RFCOMM}?
    \item Каким образом обеспечивается изоляция логических каналов в рамках одного физического соединения?
    \item Какие этапы установки RFCOMM-соединения можно наблюдать в логе \texttt{btmon}?
\end{enumerate}